\section{Introduction to Heat Exchangers}

\subsection{What is a Heat Exchanger?}
A heat exchanger is a device that facilitates heat transfer between two or more fluids (or a fluid and a solid surface) that are at different temperatures, while keeping them from mixing. They are fundamental components in many engineering applications, including power generation, refrigeration, air conditioning, chemical processing, and electronic cooling.

\subsection{Basic Principle}
The fundamental principle of a heat exchanger is the conservation of energy. Heat lost by the hot fluid (or surface) is gained by the cold fluid (or surface), assuming negligible heat losses to the surroundings. Heat transfer typically occurs across a solid wall separating the fluids.

\section{Analysis of Heat Exchangers}

\subsection{Assumptions for Heat Exchanger Analysis}
Typical assumptions made for analyzing heat exchangers operating under steady conditions include:
\begin{itemize}
    \item \textbf{Steady-State Operation:} Heat exchangers operate for long periods of time without changes in their operating conditions.
    \item \textbf{Constant Fluid Properties:} Specific heats, densities, thermal conductivities, etc., are assumed constant, particularly for fluids experiencing small temperature changes.
    \item \textbf{Negligible Kinetic and Potential Energy Changes:} Changes in kinetic and potential energy of the fluids are often negligible compared to the enthalpy changes due to heat transfer.
    \item \textbf{Adiabatic Outer Surface:} The outer surface of the heat exchanger is often assumed to be perfectly insulated, so there is no heat loss to the surrounding medium, and heat transfer occurs only between the two fluids.
    \item \textbf{No Heat Generation:} No internal heat generation within the fluid streams or the heat exchanger walls.
    \item \textbf{Axial Conduction Negligible:} Heat conduction along the tube walls in the flow direction is often assumed to be negligible.
\end{itemize}

\subsection{Energy Balances}
For a heat exchanger, the rate of heat transfer ($\Qdot$) from the hot fluid to the cold fluid can be expressed using energy balances for each fluid stream.
\begin{itemize}
    \item \textbf{For the hot fluid:} Heat released by the hot fluid is given by its mass flow rate and temperature change.
    \item \textbf{For the cold fluid:} Heat absorbed by the cold fluid is given by its mass flow rate and temperature change.
\end{itemize}
Ideally, $\Qdot_{\text{hot}} = \Qdot_{\text{cold}} = \Qdot_{\text{actual}}$.
The temperature difference between the hot and cold fluids is generally not constant throughout the heat exchanger.

\subsection{Fouling Factor ($\Rf$)}
\begin{itemize}
    \item \textbf{Definition:} Fouling refers to the accumulation of undesirable deposits on heat transfer surfaces. These deposits act as an additional thermal resistance, impeding heat transfer. The fouling factor, $\Rf$, is a measure of this thermal resistance due to the deposits.
    \item \textbf{Impact:} Fouling reduces the overall heat transfer coefficient and efficiency of the heat exchanger, requiring more surface area for the same heat transfer duty or leading to higher operating costs.
    \item \textbf{Causes:} Common causes include water hardness (scaling), corrosion, chemical reactions, and biological growth (biofouling).
\end{itemize}

\subsection{Overall Heat Transfer Coefficient ($\Uoverall$)}
\begin{itemize}
    \item \textbf{Concept:} In heat exchangers, heat transfer often involves a combination of convection (fluid to wall), conduction (through the wall), and potentially fouling. The overall heat transfer coefficient, $\Uoverall$, combines these individual resistances into a single value that characterizes the overall thermal performance.
    \item \textbf{Relation to Thermal Resistances:} The overall heat transfer coefficient is related to the sum of all individual thermal resistances (convection, conduction, fouling) in the heat transfer path.
\end{itemize}

\subsection{Log Mean Temperature Difference (LMTD) Method}
\begin{itemize}
    \item \textbf{Concept:} Since the temperature difference between the hot and cold fluids varies along the length of a heat exchanger, a special average temperature difference, the Log Mean Temperature Difference (LMTD or $\Delta T_{lm}$), is used for heat exchanger analysis.
    \item \textbf{Application:} The LMTD method is applicable for steady-flow heat exchangers where fluid properties and heat transfer coefficients are constant, and there is no phase change. It is most useful when inlet and outlet temperatures are known, allowing direct calculation of the LMTD.
    \item \textbf{Correction Factor ($F$):} For multi-pass and cross-flow heat exchangers, a correction factor $F$ is introduced to account for deviations from the pure counter-flow arrangement. This factor is typically obtained from charts based on the specific heat exchanger configuration and inlet/outlet temperatures. The modified equation is $\Qdot = \Uoverall \A F \Delta T_{lm}$.
\end{itemize}

\subsection{NTU - Effectiveness Method}
The LMTD method requires all inlet and outlet temperatures to be known. When only inlet temperatures are known and outlet temperatures are to be determined, or when phase change occurs, the NTU-Effectiveness method is often more convenient.
\subsubsection*{Heat Capacity Rates ($\Cr_h, \Cr_c$)}
\begin{itemize}
    \item The heat capacity rate for each fluid stream is defined as the product of its mass flow rate and specific heat capacity.
\end{itemize}
\subsubsection*{Maximum Possible Heat Transfer Rate ($\Qdot_{max}$)}
\begin{itemize}
    \item The maximum possible heat transfer rate in a heat exchanger is achieved when the fluid with the minimum heat capacity rate experiences the maximum possible temperature change, which is the difference between the inlet temperatures of the two fluids.
\end{itemize}
\subsubsection*{Effectiveness ($\effectiveness$)}
\begin{itemize}
    \item \textbf{Definition:} The effectiveness ($\effectiveness$) of a heat exchanger is the ratio of the actual heat transfer rate to the maximum possible heat transfer rate that could be achieved in a heat exchanger of infinite surface area.
    \item \textbf{Significance:} Effectiveness is a measure of the heat exchanger's thermal performance, ranging from 0 to 1. A higher effectiveness indicates a more efficient heat exchanger. It depends on the geometry of the heat exchanger as well as the flow arrangement.
\end{itemize}
\subsubsection*{Number of Transfer Units (NTU)}
\begin{itemize}
    \item \textbf{Definition:} The Number of Transfer Units (NTU) is a dimensionless parameter that expresses the heat transfer size of a heat exchanger. It is a measure of the heat transfer surface area, the overall heat transfer coefficient, and the minimum heat capacity rate.
    \item \textbf{Significance:} A larger NTU generally implies a larger heat exchanger and thus a higher effectiveness. However, the effectiveness increases asymptotically with NTU, meaning there are diminishing returns beyond a certain NTU value.
    \item \textbf{Relationship to Effectiveness:} For a given heat exchanger type, the effectiveness is a function of NTU and the heat capacity ratio ($\Cr = \Cmin / \Cmax$). These relationships are often presented in chart form or as analytical equations.
\end{itemize}

\section{Types of Heat Exchangers}
Heat exchangers are classified based on their construction and flow arrangement.

\subsection{Flow Arrangements}
\begin{itemize}
    \item \textbf{Parallel Flow:} Both hot and cold fluids enter at the same end and flow in the same direction. The temperature difference between the fluids decreases along the length of the heat exchanger.
    \item \textbf{Counter-Flow:} Hot and cold fluids enter at opposite ends and flow in opposite directions. This arrangement provides the largest average temperature difference, making it the most efficient for a given heat transfer surface area.
    \item \textbf{Cross-Flow:} One fluid flows perpendicular to the other fluid. This is common in compact heat exchangers and can be unmixed or mixed on one or both sides.
\end{itemize}

\subsection{Common Types of Heat Exchanger Construction}
\begin{itemize}
    \item \textbf{Double-Pipe Heat Exchanger:} Consists of two concentric pipes. One fluid flows through the inner pipe, and the other flows through the annular space between the pipes. Simple and suitable for small flow rates.
    \item \textbf{Shell-and-Tube Heat Exchanger:} Consists of a shell (large cylindrical vessel) with a bundle of tubes inside. One fluid flows through the tubes (tube side), and the other fluid flows outside the tubes, through the shell (shell side). Baffles are often used in the shell to promote turbulence and improve heat transfer. These are widely used in industrial applications due to their robustness and high heat transfer capacity.
    \item \textbf{Plate and Frame Heat Exchanger:} Consists of a series of thin, corrugated plates pressed together, creating channels for the fluids. They offer high heat transfer coefficients and are compact, easy to clean, and allow for easy expansion or reduction of capacity.
\end{itemize}

\section{Heat Exchanger Mechanism}
Heat transfer in a heat exchanger involves a series of resistances:
\begin{itemize}
    \item 1.  Convection from the hot fluid to the inner surface of the tube.
    \item 2.  Conduction through any fouling layer on the inner surface.
    \item 3.  Conduction through the tube wall.
    \item 4.  Conduction through any fouling layer on the outer surface.
    \item 5.  Convection from the outer surface of the tube to the cold fluid.
\end{itemize}


These resistances are typically in series, similar to an electrical resistance network.

 % Start formulas on a new page for clarity
\section{Formulas Summary for Heat Exchangers}

\subsection{Energy Balances}
\subsubsection*{Heat Transfer Rate based on Hot Fluid}
\begin{equation}
    \dot{Q} = \dot{m}_h C_{p,h} (T_{h,in} - T_{h,out})
\end{equation}
where: \\
$\Qdot$ = Rate of heat transfer (W) \\
$\mdot_h$ = Mass flow rate of the hot fluid (kg/s) \\
$Cp_{h}$ = Specific heat capacity of the hot fluid (J/kg$\cdot$K) \\
$T_{h,in}$ = Inlet temperature of the hot fluid (K or $^\circ$C) \\
$T_{h,out}$ = Outlet temperature of the hot fluid (K or $^\circ$C)

\subsubsection*{Heat Transfer Rate based on Cold Fluid}
\begin{equation}
    \Qdot = \mdot_c C_{p,c} (\T_{c,out} - \T_{c,in})
\end{equation}
where: \\
$\mdot_c$ = Mass flow rate of the cold fluid (kg/s) \\
$C_{p,c}$ = Specific heat capacity of the cold fluid (J/kg$\cdot$K) \\
$\T_{c,out}$ = Outlet temperature of the cold fluid (K or $^\circ$C) \\
$\T_{c,in}$ = Inlet temperature of the cold fluid (K or $^\circ$C)

\subsection{Fouling Factor ($\Rf$)}
Fouling resistance for inner surface: $R_{f,i}$ (m$^2$$\cdot$K/W) \\
Fouling resistance for outer surface: $R_{f,o}$ (m$^2$$\cdot$K/W)

\subsection{Overall Heat Transfer Coefficient ($\Uoverall$)}
\subsubsection*{General Relation}
\begin{equation}
    \Qdot = \Uoverall \A \Delta T_{\text{avg}}
\end{equation}
where: \\
$\Uoverall$ = Overall heat transfer coefficient (W/m$^2$$\cdot$K) \\
$\A$ = Total heat transfer surface area (m$^2$) \\
$\Delta T_{\text{avg}}$ = Average temperature difference (e.g., LMTD) (K or $^\circ$C)

\subsubsection*{Overall Thermal Resistance for a Cylindrical Tube (based on $\A_i$ or $\A_o$)}
\begin{equation}
    \frac{1}{\Uoverall_i \A_i} = R_{\text{total}} = \frac{1}{\hC_i \A_i} + \frac{R_{f,i}}{\A_i} + \frac{\ln(D_o/D_i)}{2\pi L \kth} + \frac{R_{f,o}}{\A_o} + \frac{1}{\hC_o \A_o}
\end{equation}
or, expressed per unit area for $\Uoverall_i$:
\begin{equation}
    \frac{1}{\Uoverall_i} = \frac{1}{\hC_i} + R_{f,i} + \frac{A_i \ln(D_o/D_i)}{2\pi L \kth} + \frac{A_i}{A_o} R_{f,o} + \frac{A_i}{A_o} \frac{1}{\hC_o}
\end{equation}
\begin{equation}
    \frac{1}{\Uoverall_i} = \frac{1}{\hC_i} + R_{f,i} + \frac{D_i \ln(D_o/D_i)}{2\kth} + \frac{D_i}{D_o} R_{f,o} + \frac{D_i}{D_o} \frac{1}{\hC_o}
\end{equation}
where: \\
$\Uoverall_i$ = Overall heat transfer coefficient based on inner surface area (W/m$^2$$\cdot$K) \\
$\A_i$ = Inner surface area (m$^2$) \\
$\A_o$ = Outer surface area (m$^2$) \\
$D_i, D_o$ = Inner and outer diameters of the tube (m) \\
$L$ = Length of the tube (m) \\
$\kth$ = Thermal conductivity of the tube wall (W/m$\cdot$K) \\
$\hC_i, \hC_o$ = Convection heat transfer coefficients on inner and outer surfaces (W/m$^2$$\cdot$K)

\subsection{Log Mean Temperature Difference (LMTD)}
\subsubsection*{For Parallel Flow}
\begin{align}
    \Delta T_1 &= \T_{h,in} - \T_{c,in} \\
    \Delta T_2 &= \T_{h,out} - \T_{c,out} \\
    \Delta T_{lm,parallel} &= \frac{\Delta T_1 - \Delta T_2}{\ln(\Delta T_1 / \Delta T_2)}
\end{align}

\subsubsection*{For Counter-Flow}
\begin{align}
    \Delta T_1 &= \T_{h,in} - \T_{c,out} \\
    \Delta T_2 &= \T_{h,out} - \T_{c,in} \\
    \Delta T_{lm,counter} &= \frac{\Delta T_1 - \Delta T_2}{\ln(\Delta T_1 / \Delta T_2)}
\end{align}
where: \\
$\Delta T_{lm}$ = Log Mean Temperature Difference (K or $^\circ$C) \\
$\Delta T_1, \Delta T_2$ = Temperature differences at the two ends of the heat exchanger (K or $^\circ$C) \\
$\T_{h,in}, \T_{h,out}$ = Inlet and outlet temperatures of the hot fluid (K or $^\circ$C) \\
$\T_{c,in}, \T_{c,out}$ = Inlet and outlet temperatures of the cold fluid (K or $^\circ$C)

\subsection{NTU - Effectiveness Method}
\subsubsection*{Heat Capacity Rates}

\begin{align}
    C_h &= \mdot_h Cp_{h} && \text{(Heat capacity rate of hot fluid, W/K)} \\
    C_c &= \mdot_c C_{p,c} && \text{(Heat capacity rate of cold fluid, W/K)}
\end{align}

\begin{align}
    \Cmin &= \min(C_h, C_c) && \text{(Minimum heat capacity rate, W/K)} \\
    \Cmax &= \max(C_h, C_c) && \text{(Maximum heat capacity rate, W/K)}
\end{align}
\begin{equation}
    \Cr = \frac{\Cmin}{\Cmax} \quad \text{(Heat capacity ratio, dimensionless)}
\end{equation}

\subsubsection*{Maximum Possible Heat Transfer Rate}
\begin{equation}
    \Qdot_{max} = \Cmin (\T_{h,in} - \T_{c,in})
\end{equation}
where: \\
$\Qdot_{max}$ = Maximum possible heat transfer rate (W) \\
$\T_{h,in}$ = Inlet temperature of the hot fluid (K or $^\circ$C) \\
$\T_{c,in}$ = Inlet temperature of the cold fluid (K or $^\circ$C)

\subsubsection*{Actual Heat Transfer Rate (from Effectiveness)}
\begin{equation}
    \Qdot = \effectiveness \Qdot_{max}
\end{equation}
where: \\
$\effectiveness$ = Heat exchanger effectiveness (dimensionless)

\subsubsection*{Effectiveness ($\effectiveness$)}
\begin{equation}
    \effectiveness = \frac{\text{Actual heat transfer rate}}{\text{Maximum possible heat transfer rate}} = \frac{\Qdot}{\Qdot_{max}}
\end{equation}
For hot fluid: $\effectiveness = \frac{\mdot_h C_{p_h} (\T_{h,in} - \T_{h,out})}{\Cmin (\T_{h,in} - \T_{c,in})}$ \\
For cold fluid: $\effectiveness = \frac{\mdot_c C_{p,c} (\T_{c,out} - \T_{c,in})}{\Cmin (\T_{h,in} - \T_{c,in})}$

\subsubsection*{Number of Transfer Units (NTU)}
\begin{equation}
    \NTU = \frac{\Uoverall \A}{\Cmin}
\end{equation}
where: \\
$\NTU$ = Number of Transfer Units (dimensionless) \\
$\Uoverall$ = Overall heat transfer coefficient (W/m$^2$$\cdot$K) \\
$\A$ = Heat transfer surface area (m$^2$) \\
$\Cmin$ = Minimum heat capacity rate (W/K)

\subsubsection*{Effectiveness-NTU Relations for various Flow Arrangements}
The effectiveness $\effectiveness$ is a function of $\NTU$ and $\Cr$: $\effectiveness = f(\NTU, \Cr)$

\begin{itemize}
    \item \textbf{Parallel Flow:}
    \begin{equation}
        \effectiveness = \frac{1 - \exp[-\NTU(1 + \Cr)]}{1 + \Cr}
    \end{equation}

    \item \textbf{Counter-Flow:}
    \begin{itemize}
        \item If $\Cr < 1$:
        \begin{equation}
            \effectiveness = \frac{1 - \exp[-\NTU(1 - \Cr)]}{1 - \Cr \exp[-\NTU(1 - \Cr)]}
        \end{equation}
        \item If $\Cr = 1$:
        \begin{equation}
            \effectiveness = \frac{\NTU}{1 + \NTU}
        \end{equation}
    \end{itemize}

    \item \textbf{Shell and Tube (1 shell pass, 2, 4, 6, ... tube passes):}
    \begin{equation}
        \effectiveness = 2 \left\{ 1 + \Cr + (1 + \Cr^2)^{1/2} \frac{1 + \exp[-\NTU(1 + \Cr^2)^{1/2}]}{1 - \exp[-\NTU(1 + \Cr^2)^{1/2}]} \right\}^{-1}
    \end{equation}

    \item \textbf{Cross-Flow (Single-pass):}
    \begin{itemize}
        \item \textbf{Both Fluids Unmixed:}
        \begin{equation}
            \effectiveness = 1 - \exp \left[ \frac{1}{\Cr} \NTU^{0.22} \left( \exp(-\Cr \NTU^{0.78}) - 1 \right) \right]
        \end{equation}
        \item \textbf{$\Cmax$ Mixed, $\Cmin$ Unmixed:}
        \begin{equation}
            \effectiveness = \frac{1}{\Cr} [1 - \exp(-\Cr(1 - e^{-\NTU}))]
        \end{equation}
        \item \textbf{$\Cmin$ Mixed, $\Cmax$ Unmixed:}
        \begin{equation}
            \effectiveness = 1 - \exp \left( -\frac{1}{\Cr} [1 - e^{-\Cr \NTU}] \right)
        \end{equation}
    \end{itemize}
\end{itemize}
